\documentclass[10pt]{article}
\usepackage[a4paper,margin=5mm]{geometry}
\usepackage{amsmath}
\usepackage{fontspec}
\usepackage{unicode-math}
\setmainfont{NewCM10}[Extension=.otf,UprightFont=*-Book,BoldFont=*-Bold,ItalicFont=*-BookItalic,BoldItalicFont=*-BoldItalic]
\setsansfont{NewCMSans10}[Extension=.otf,UprightFont=*-Book,BoldFont=*-Bold,ItalicFont=*-BookOblique,BoldItalicFont=*-BoldOblique]
\setmathfont{NewCMMath-Book.otf}
\usepackage{xcolor}
\definecolor{FAUblue}{HTML}{04316A}
\definecolor{Paperblue}{HTML}{EEF2F7}
\usepackage{multicol,enumitem,tabularx,array,microtype,tikz,needspace}
\usepackage[hidelinks]{hyperref}
\usepackage{draftwatermark}
\SetWatermarkText{Made by Thanate Amornjiraporn (Focus)}
\SetWatermarkColor{black!5}\SetWatermarkFontSize{30pt}\SetWatermarkAngle{45}
\pagestyle{empty}
\setlength{\parindent}{0pt}\setlength{\parskip}{0.4pt}
\setlength{\columnsep}{2.6mm}\setlength{\columnseprule}{0.25pt}\setlength{\multicolsep}{0.5pt}
\setlist{nosep,leftmargin=0.85em,labelsep=0.2em,label={\color{FAUblue}\textbullet}}
\setlist[enumerate]{nosep,leftmargin=1.05em,labelsep=0.2em,label={\color{FAUblue}\arabic*.}}
\setlength{\abovedisplayskip}{1pt}\setlength{\belowdisplayskip}{1pt}
\setlength{\lineskiplimit}{0pt}\setlength{\lineskip}{0.4pt}
\newcommand\sheethead[1]{\par\needspace{3\baselineskip}\vspace{1.6pt}{\color{FAUblue}\sffamily\bfseries #1}\par\nopagebreak\vspace{0.3pt}}
\newcommand\sheetf[1]{\par\hspace*{0.5em}#1\par}
\begin{document}
{\sffamily\bfseries\color{FAUblue}\small Side 1 \ensuremath{\cdot} Rules from Exercises 1–4}\hfill{\fontsize{5}{6}\selectfont\color{black!65}Appendix B: Exam sheet, Exercises 1–4 · Mathematics for Engineers I \ensuremath{\cdot} Module 27651 \ensuremath{\cdot} Appendix B}\par\vspace{0.6pt}{\color{FAUblue}\hrule height 0.5pt}\vspace{1.2pt}
\begin{multicols*}{4}\fontsize{7.90}{9.24}\selectfont\raggedright
{\sheethead{Statements and connectives}\par A \textbf{statement} is a sentence that is either true (T) or false (F), never both. “The formula is beautiful” is not one.\par\par{\setlength{\tabcolsep}{2.4pt}\renewcommand{\arraystretch}{0.92}\begin{tabularx}{\linewidth}{@{}lllll>{\raggedright\arraybackslash}X@{}}
\textbf{\textit{A} \textit{B}} & \textbf{\ensuremath{\neg}\textit{A}} & \textbf{\textit{A}\ensuremath{\land}\textit{B}} & \textbf{\textit{A}\ensuremath{\lor}\textit{B}} & \textbf{\textit{A}\ensuremath{\Rightarrow}\textit{B}} & \textbf{\textit{A}\ensuremath{\Leftrightarrow}\textit{B}} \\
T T & F & T & T & T & T \\
T F & F & F & T & F & F \\
F T & T & F & T & T & F \\
F F & T & F & F & T & T \\
\end{tabularx}}\par\par{\color{black!72}\ensuremath{\lor} is inclusive (both may be true). \textit{A} \ensuremath{\Rightarrow} \textit{B} is false only in the row T \ensuremath{\Rightarrow} F, so a false premise makes it true. \textbf{Tautology}: T in every row; \textbf{contradiction}: F in every row. \textit{n} letters give 2\textsuperscript{\textit{n}} rows.}\par}
{\sheethead{Equivalences to know}\begin{itemize}
\item \ensuremath{\neg}\ensuremath{\neg}\textit{A} \ensuremath{\Leftrightarrow} \textit{A}
\item \textit{A} \ensuremath{\Rightarrow} \textit{B} \ensuremath{\Leftrightarrow} \ensuremath{\neg}\textit{A} \ensuremath{\lor} \textit{B} \ensuremath{\Leftrightarrow} \ensuremath{\neg}\textit{B} \ensuremath{\Rightarrow} \ensuremath{\neg}\textit{A} (contrapositive)
\item \ensuremath{\neg}(\textit{A} \ensuremath{\Rightarrow} \textit{B}) \ensuremath{\Leftrightarrow} \textit{A} \ensuremath{\land} \ensuremath{\neg}\textit{B}
\item \textit{A} \ensuremath{\Leftrightarrow} \textit{B} \ensuremath{\Leftrightarrow} (\textit{A} \ensuremath{\Rightarrow} \textit{B}) \ensuremath{\land} (\textit{B} \ensuremath{\Rightarrow} \textit{A})
\item De Morgan: \ensuremath{\neg}(\textit{A} \ensuremath{\land} \textit{B}) \ensuremath{\Leftrightarrow} \ensuremath{\neg}\textit{A} \ensuremath{\lor} \ensuremath{\neg}\textit{B}; \ensuremath{\neg}(\textit{A} \ensuremath{\lor} \textit{B}) \ensuremath{\Leftrightarrow} \ensuremath{\neg}\textit{A} \ensuremath{\land} \ensuremath{\neg}\textit{B}
\item \textit{A} \ensuremath{\land} (\textit{B} \ensuremath{\lor} \textit{C}) \ensuremath{\Leftrightarrow} (\textit{A} \ensuremath{\land} \textit{B}) \ensuremath{\lor} (\textit{A} \ensuremath{\land} \textit{C}); swap \ensuremath{\land}, \ensuremath{\lor} for the twin
\end{itemize}\par{\color{black!72}The converse \textit{B} \ensuremath{\Rightarrow} \textit{A} is \textbf{not} equivalent to \textit{A} \ensuremath{\Rightarrow} \textit{B}. Prove an equivalence: equal truth columns in every row. Show “not a tautology”: one row giving F.}\par}
{\sheethead{Words to logic}\par{\setlength{\tabcolsep}{2.4pt}\renewcommand{\arraystretch}{0.92}\begin{tabularx}{\linewidth}{@{}l>{\raggedright\arraybackslash}X@{}}
\textit{A} if \textit{B} & \textit{B} \ensuremath{\Rightarrow} \textit{A} (\textit{B} is sufficient) \\
\textit{A} only if \textit{B} & \textit{A} \ensuremath{\Rightarrow} \textit{B} (\textit{B} is necessary) \\
\textit{A} if and only if \textit{B} & \textit{A} \ensuremath{\Leftrightarrow} \textit{B} \\
\textit{A} unless \textit{B} & \ensuremath{\neg}\textit{B} \ensuremath{\Rightarrow} \textit{A}, same as \textit{A} \ensuremath{\lor} \textit{B} \\
not \textit{A} unless \textit{B} & \ensuremath{\neg}\textit{B} \ensuremath{\Rightarrow} \ensuremath{\neg}\textit{A}, same as \textit{A} \ensuremath{\Rightarrow} \textit{B} \\
neither \textit{A} nor \textit{B} & \ensuremath{\neg}\textit{A} \ensuremath{\land} \ensuremath{\neg}\textit{B} \\
“but” & \ensuremath{\land} (and) \\
\end{tabularx}}\par}
{\sheethead{Quantifiers}\begin{itemize}
\item \ensuremath{\forall}\textit{x} \ensuremath{\in} \textit{M}: P(\textit{x}) — for \textbf{all} \textit{x}
\item \ensuremath{\exists}\textit{x} \ensuremath{\in} \textit{M}: P(\textit{x}) — for \textbf{at least one} \textit{x}
\item \ensuremath{\exists}!\textit{x} — exactly one
\end{itemize}\sheetf{\ensuremath{\neg \forall {x}\, \mathrm{P}({x})\, \Leftrightarrow \, \exists {x}\, \neg \mathrm{P}({x})}}\sheetf{\ensuremath{\neg \exists {x}\, \mathrm{P}({x})\, \Leftrightarrow \, \forall {x}\, \neg \mathrm{P}({x})}}\par{\color{black!72}Negate: swap every \ensuremath{\forall} \ensuremath{\leftrightarrow} \ensuremath{\exists} and negate the end: \ensuremath{\neg}(\ensuremath{\forall}\textit{x} \ensuremath{\exists}\textit{y}: \textit{x} \ensuremath{<} \textit{y}) \ensuremath{\Leftrightarrow} \ensuremath{\exists}\textit{x} \ensuremath{\forall}\textit{y}: \textit{x} \ensuremath{\geq} \textit{y}. Order matters: \ensuremath{\forall}\textit{y} \ensuremath{\exists}\textit{x}: \textit{x} \ensuremath{<} \textit{y} is true in \ensuremath{\mathbb{Z}}, \ensuremath{\exists}\textit{x} \ensuremath{\forall}\textit{y}: \textit{x} \ensuremath{<} \textit{y} is false everywhere (take \textit{y} = \textit{x}). One counterexample kills a \ensuremath{\forall}; one example proves an \ensuremath{\exists}.}\par}
{\sheethead{Proof patterns}\begin{itemize}
\item Direct: assume \textit{A}, derive \textit{B}.
\item Contrapositive: assume \ensuremath{\neg}\textit{B}, derive \ensuremath{\neg}\textit{A}.
\item Contradiction: assume the claim is false, reach something false.
\item \textit{A} \ensuremath{\Leftrightarrow} \textit{B}: prove both directions.
\item Set equality: two inclusions, or a chain of \ensuremath{\Leftrightarrow} on “\textit{x} \ensuremath{\in} …”.
\end{itemize}}
{\sheethead{Sets: notation}\begin{itemize}
\item \textit{a} \ensuremath{\in} \textit{A}: \textit{a} is an element of \textit{A}; \ensuremath{\emptyset}: the empty set; |\textit{A}|: number of elements
\item \textit{A} \ensuremath{\subset} \textit{B}: every element of \textit{A} lies in \textit{B} (\ensuremath{\emptyset} \ensuremath{\subset} \textit{A} and \textit{A} \ensuremath{\subset} \textit{A} always)
\item \textit{A} = \textit{B} \ensuremath{\Leftrightarrow} \textit{A} \ensuremath{\subset} \textit{B} and \textit{B} \ensuremath{\subset} \textit{A}
\item [\textit{a}, \textit{b}] closed (ends in), (\textit{a}, \textit{b}) open (ends out), [\textit{a}, \textit{b}) half-open
\item \ensuremath{\mathbb{N}} = \{1, 2, 3, …\}, \ensuremath{\mathbb{N}}\textsubscript{0} = \ensuremath{\mathbb{N}} \ensuremath{\cup} \{0\}; \ensuremath{\mathbb{N}} \ensuremath{\subsetneq} \ensuremath{\mathbb{Z}} \ensuremath{\subsetneq} \ensuremath{\mathbb{Q}} \ensuremath{\subsetneq} \ensuremath{\mathbb{R}}
\end{itemize}}
{\sheethead{Set operations}\begin{itemize}
\item \textit{A} \ensuremath{\cup} \textit{B}: in \textit{A} \textbf{or} \textit{B}; \textit{A} \ensuremath{\cap} \textit{B}: in \textbf{both}
\item \textit{A} \textbackslash{} \textit{B}: in \textit{A} and not in \textit{B} (= \textit{A} \ensuremath{\cap} \textit{B}\textsuperscript{c})
\item \textit{A}\textsuperscript{c} = \textit{X} \textbackslash{} \textit{A}: always say what \textit{X} is
\item \textit{A} \ensuremath{\times} \textit{B} = \{(\textit{a}, \textit{b}) : \textit{a} \ensuremath{\in} \textit{A}, \textit{b} \ensuremath{\in} \textit{B}\}: ordered pairs; |\textit{A} \ensuremath{\times} \textit{B}| = |\textit{A}| \ensuremath{\cdot} |\textit{B}|; \textit{A} \ensuremath{\times} \ensuremath{\emptyset} = \ensuremath{\emptyset}
\item \ensuremath{\mathcal{P}}(\textit{A}): all subsets, \ensuremath{\emptyset} and \textit{A} included; |\ensuremath{\mathcal{P}}(\textit{A})| = 2\textsuperscript{|\textit{A}|}; \textit{X} \ensuremath{\in} \ensuremath{\mathcal{P}}(\textit{A}) \ensuremath{\Leftrightarrow} \textit{X} \ensuremath{\subset} \textit{A}
\end{itemize}\par{\color{black!72}\{1, 2\} \ensuremath{\notin} \ensuremath{\mathcal{P}}(\{1, 2\} \ensuremath{\times} \{1\}): the elements of that product are pairs, not numbers.}\par}
{\sheethead{Set laws}\begin{itemize}
\item (\textit{A} \ensuremath{\cup} \textit{B})\textsuperscript{c} = \textit{A}\textsuperscript{c} \ensuremath{\cap} \textit{B}\textsuperscript{c}; (\textit{A} \ensuremath{\cap} \textit{B})\textsuperscript{c} = \textit{A}\textsuperscript{c} \ensuremath{\cup} \textit{B}\textsuperscript{c}
\item \textit{A} \ensuremath{\cap} (\textit{B} \ensuremath{\cup} \textit{C}) = (\textit{A} \ensuremath{\cap} \textit{B}) \ensuremath{\cup} (\textit{A} \ensuremath{\cap} \textit{C}), and with \ensuremath{\cap}, \ensuremath{\cup} swapped
\item \textit{A} \textbackslash{} (\textit{B} \ensuremath{\cup} \textit{C}) = (\textit{A} \textbackslash{} \textit{B}) \ensuremath{\cap} (\textit{A} \textbackslash{} \textit{C}); \textit{A} \textbackslash{} (\textit{B} \ensuremath{\cap} \textit{C}) = (\textit{A} \textbackslash{} \textit{B}) \ensuremath{\cup} (\textit{A} \textbackslash{} \textit{C})
\item \textit{A} \textbackslash{} (\textit{B} \textbackslash{} \textit{C}) = (\textit{A} \textbackslash{} \textit{B}) \ensuremath{\cup} (\textit{A} \ensuremath{\cap} \textit{C})
\end{itemize}\par{\color{black!72}Logic \ensuremath{\leftrightarrow} sets: \ensuremath{\land} \ensuremath{\leftrightarrow} \ensuremath{\cap}, \ensuremath{\lor} \ensuremath{\leftrightarrow} \ensuremath{\cup}, \ensuremath{\neg} \ensuremath{\leftrightarrow} complement, \ensuremath{\Rightarrow} \ensuremath{\leftrightarrow} \ensuremath{\subset}. Intervals: check each endpoint; (1, 3)\textsuperscript{c} in \ensuremath{\mathbb{R}} is (\ensuremath{-}\ensuremath{\infty}, 1] \ensuremath{\cup} [3, \ensuremath{\infty}).}\par}
{\sheethead{Order axioms: consequences}\begin{itemize}
\item \textit{a} \ensuremath{\cdot} 0 = 0 (from \textit{a} \ensuremath{\cdot} 0 = \textit{a}(0 + 0) = \textit{a} \ensuremath{\cdot} 0 + \textit{a} \ensuremath{\cdot} 0)
\item \ensuremath{-}\textit{a} = (\ensuremath{-}1) \ensuremath{\cdot} \textit{a}; (\ensuremath{-}\textit{a})(\ensuremath{-}\textit{b}) = \textit{a}\textit{b}
\item \textit{a}\textit{b} = 0 \ensuremath{\Rightarrow} \textit{a} = 0 or \textit{b} = 0
\item \textit{a} \ensuremath{\neq} 0 \ensuremath{\Rightarrow} \textit{a}\textsuperscript{2} \ensuremath{>} 0; hence 1 = 1\textsuperscript{2} \ensuremath{>} 0
\item \textit{a} \ensuremath{<} \textit{b} \ensuremath{\Rightarrow} \textit{a} + \textit{c} \ensuremath{<} \textit{b} + \textit{c}
\item \textit{a} \ensuremath{<} \textit{b} and \textit{c} \ensuremath{>} 0 \ensuremath{\Rightarrow} \textit{a}\textit{c} \ensuremath{<} \textit{b}\textit{c}; \textit{c} \ensuremath{<} 0 \ensuremath{\Rightarrow} \textit{a}\textit{c} \ensuremath{>} \textit{b}\textit{c} (flip!)
\item 0 \ensuremath{<} \textit{a} \ensuremath{<} \textit{b} \ensuremath{\Rightarrow} {\ensuremath{\dfrac{1}{{a}}}} \ensuremath{>} {\ensuremath{\dfrac{1}{{b}}}}
\end{itemize}\par{\color{black!72}In an axiom proof cite one axiom per step: commutative, associative, distributive, 0 and 1, inverses, trichotomy, transitivity, + and positive \ensuremath{\cdot} keep order.}\par}
{\sheethead{Bounds, supremum, infimum}\begin{itemize}
\item Upper bound \textit{s}: \textit{x} \ensuremath{\leq} \textit{s} for every \textit{x} \ensuremath{\in} \textit{M}
\item sup \textit{M}: the \textbf{least} upper bound; inf \textit{M}: the greatest lower bound
\item max \textit{M}: a sup that lies in \textit{M} (same for min, inf)
\item Completeness: every non-empty \textit{M} \ensuremath{\subset} \ensuremath{\mathbb{R}} bounded above has a sup in \ensuremath{\mathbb{R}}
\end{itemize}\par{\color{black!72}Show \textit{s} = sup \textit{M}: (1) \textit{x} \ensuremath{\leq} \textit{s} for all \textit{x} \ensuremath{\in} \textit{M}; (2) for every ε \ensuremath{>} 0 some \textit{x} \ensuremath{\in} \textit{M} has \textit{x} \ensuremath{>} \textit{s} \ensuremath{-} ε.}\par\par{\setlength{\tabcolsep}{2.4pt}\renewcommand{\arraystretch}{0.92}\begin{tabularx}{\linewidth}{@{}llll>{\raggedright\arraybackslash}X@{}}
\textbf{Set} & \textbf{sup} & \textbf{max} & \textbf{inf} & \textbf{min} \\
(0, 1] & 1 & 1 & 0 & none \\
\{{\ensuremath{\dfrac{1}{{n}}}} : \textit{n} \ensuremath{\in} \ensuremath{\mathbb{N}}\} & 1 & 1 & 0 & none \\
\{(\ensuremath{-}1)\textsuperscript{\textit{n}}(1 + {\ensuremath{\dfrac{1}{{n}}}})\} & {\ensuremath{\dfrac{3}{2}}} & {\ensuremath{\dfrac{3}{2}}} & \ensuremath{-}2 & \ensuremath{-}2 \\
\end{tabularx}}\par\par{\color{black!72}\ensuremath{\mathbb{Q}} is not complete: \{\textit{x} \ensuremath{\in} \ensuremath{\mathbb{Q}} : \textit{x}\textsuperscript{2} \ensuremath{<} 2\} has rational upper bounds but its sup \ensuremath{\sqrt{2}} is not rational.}\par}
{\sheethead{Absolute value}\sheetf{\ensuremath{|{x}|\, =\, {x}\, \text{if}\, {x}\, \geq \, 0,\, -{x}\, \text{if}\, {x}\, <\, 0;\, |{x}|\, =\, \sqrt{{x}^{2}}}}\begin{itemize}
\item |\textit{x} \ensuremath{-} \textit{a}| is the distance from \textit{x} to \textit{a}
\item |\textit{x}\textit{y}| = |\textit{x}| |\textit{y}|; |\textit{x}| \ensuremath{\leq} \textit{a} \ensuremath{\Leftrightarrow} \ensuremath{-}\textit{a} \ensuremath{\leq} \textit{x} \ensuremath{\leq} \textit{a}; |\textit{x}| \ensuremath{\geq} \textit{a} \ensuremath{\Leftrightarrow} \textit{x} \ensuremath{\leq} \ensuremath{-}\textit{a} or \textit{x} \ensuremath{\geq} \textit{a}
\item Triangle: |\textit{x} + \textit{y}| \ensuremath{\leq} |\textit{x}| + |\textit{y}|; ||\textit{x}| \ensuremath{-} |\textit{y}|| \ensuremath{\leq} |\textit{x} \ensuremath{-} \textit{y}|
\end{itemize}}
{\sheethead{Solving with | | and roots}\begin{enumerate}
\item Find where each |…| is zero; these points cut \ensuremath{\mathbb{R}} into pieces.
\item On each piece drop the bars with the right sign (|\textit{x} \ensuremath{-} 3| = 3 \ensuremath{-} \textit{x} for \textit{x} \ensuremath{<} 3).
\item Solve, keep only solutions inside that piece, unite the pieces.
\end{enumerate}\par{\color{black!72}Roots: radicand \ensuremath{\geq} 0 first. For a root \ensuremath{\leq} \textit{R} also need \textit{R} \ensuremath{\geq} 0. Square only when both sides are \ensuremath{\geq} 0, then check the ends. Regions: |\textit{x}| + |\textit{y}| = \textit{r} is a diamond with corners (\ensuremath{\pm}\textit{r}, 0), (0, \ensuremath{\pm}\textit{r}); |\textit{x} \ensuremath{-} \textit{a}| + |\textit{y} \ensuremath{-} \textit{b}| moves its centre to (\textit{a}, \textit{b}); \ensuremath{<} inside, \ensuremath{\leq} with the edge.}\par}
{\sheethead{Sums and products}\sheetf{\ensuremath{\sum_{{k}={m}}^{{n}} \, {a}_{{k}}\, =\, {a}_{{m}}\, +\, …\, +\, {a}_{{n}}\, ({n}\, -\, {m}\, +\, 1\, \text{terms})}}\begin{itemize}
\item \ensuremath{\textstyle\sum}(\textit{a\textsubscript{k}} + \textit{b\textsubscript{k}}) = \ensuremath{\textstyle\sum}\textit{a\textsubscript{k}} + \ensuremath{\textstyle\sum}\textit{b\textsubscript{k}}; \ensuremath{\textstyle\sum} \textit{c}\textit{a\textsubscript{k}} = \textit{c} \ensuremath{\textstyle\sum}\textit{a\textsubscript{k}}; $\sum_{{k}=1}^{{n}} $ \textit{c} = \textit{n}\textit{c}
\item Shift: $\sum_{{k}=1}^{{n}} $ \textit{a}\textsubscript{\textit{k}} = $\sum_{{j}=0}^{{n}-1} $ \textit{a}\textsubscript{\textit{j}+1}
\item Telescoping: \ensuremath{\textstyle\sum}(\textit{a}\textsubscript{\textit{k}+1} \ensuremath{-} \textit{a}\textsubscript{\textit{k}}) = \textit{a}\textsubscript{\textit{n}+1} \ensuremath{-} \textit{a}\textsubscript{1}; \ensuremath{\textstyle\prod} {\ensuremath{\dfrac{{a}_{{k}+1}}{{a}_{{k}}}}} = {\ensuremath{\dfrac{{a}_{{n}+1}}{{a}_{1}}}}
\item \ensuremath{\textstyle\prod} \textit{c} = \textit{c}\textsuperscript{\textit{n}} (\textit{n} factors); empty sum = 0, empty product = 1
\end{itemize}}
{\sheethead{Standard sums}\sheetf{\ensuremath{\sum_{{k}=0}^{{n}} \, {q}^{{k}}\, =\, \dfrac{1\, -\, {q}^{{n}+1}}{1\, -\, {q}}\, ({q}\, ≠\, 1)}}\sheetf{\ensuremath{\sum_{{k}=1}^{{n}} \, {k}\, =\, \dfrac{{n}({n}\, +\, 1)}{2}}}\sheetf{\ensuremath{\sum_{{k}=1}^{{n}} \, {k}^{2}\, =\, \dfrac{{n}({n}\, +\, 1)(2{n}\, +\, 1)}{6}}}\sheetf{\ensuremath{\sum_{{k}=1}^{{n}} \, {k}^{3}\, =\, \left(\dfrac{{n}({n}\, +\, 1)}{2}\right)^{2}}}\sheetf{\ensuremath{\sum_{{k}=1}^{{n}} \, (2{k}\, -\, 1)\, =\, {n}^{2}\, (\text{odd}\, \text{numbers})}}}
{\sheethead{Factorial, binomial coefficient}\sheetf{\ensuremath{{n}!\, =\, 1\, \cdot \, 2\, \cdot \, …\, \cdot \, {n},\, 0!\, =\, 1}}\sheetf{\ensuremath{\mathrm{C}({n},\, {k})\, =\, \dfrac{{n}!}{{k}!\, ({n}\, -\, {k})!}\, =\, \dfrac{{n}({n}\, -\, 1)…({n}\, -\, {k}\, +\, 1)}{{k}!}}}\begin{itemize}
\item C(\textit{n}, 0) = C(\textit{n}, \textit{n}) = 1; C(\textit{n}, \textit{k}) = C(\textit{n}, \textit{n} \ensuremath{-} \textit{k})
\item Pascal: C(\textit{n}, \textit{k} \ensuremath{-} 1) + C(\textit{n}, \textit{k}) = C(\textit{n} + 1, \textit{k})
\item Binomial theorem: (\textit{a} + \textit{b})\textsuperscript{\textit{n}} = $\sum_{{k}=0}^{{n}} $ C(\textit{n}, \textit{k}) \textit{a}\textsuperscript{\textit{n}\ensuremath{-}\textit{k}}\textit{b}\textsuperscript{\textit{k}}
\item \ensuremath{\textstyle\sum} C(\textit{n}, \textit{k}) = 2\textsuperscript{\textit{n}}; \ensuremath{\textstyle\sum} (\ensuremath{-}1)\textsuperscript{\textit{k}} C(\textit{n}, \textit{k}) = 0 (\textit{n} \ensuremath{\geq} 1)
\end{itemize}\par{\color{black!72}Rows: 1 \ensuremath{\cdot} 1 1 \ensuremath{\cdot} 1 2 1 \ensuremath{\cdot} 1 3 3 1 \ensuremath{\cdot} 1 4 6 4 1 \ensuremath{\cdot} 1 5 10 10 5 1 \ensuremath{\cdot} 1 6 15 20 15 6 1. Diagonal sums give Fibonacci: \textit{F}\textsubscript{1} = \textit{F}\textsubscript{2} = 1, \textit{F}\textsubscript{\textit{n}+2} = \textit{F}\textsubscript{\textit{n}+1} + \textit{F}\textsubscript{\textit{n}}.}\par}
{\sheethead{Counting}\par{\setlength{\tabcolsep}{2.4pt}\renewcommand{\arraystretch}{0.92}\begin{tabularx}{\linewidth}{@{}l>{\raggedright\arraybackslash}X@{}}
\textbf{Choose \textit{k} of \textit{n}} & \textbf{Count} \\
order ignored & C(\textit{n}, \textit{k}) \\
order matters, no repeats & {\ensuremath{\dfrac{{n}!}{({n}\, -\, {k})!}}} \\
order matters, repeats allowed & \textit{n}\textsuperscript{\textit{k}} \\
arrange all \textit{n} & \textit{n}! \\
\end{tabularx}}\par\begin{itemize}
\item Each pair once (handshakes, clinks): C(\textit{n}, 2) = {\ensuremath{\dfrac{{n}({n}\, -\, 1)}{2}}}; 66 clinks \ensuremath{\Rightarrow} \textit{n} = 12
\item Each ordered pair (speaker \ensuremath{\to} listener): \textit{n}(\textit{n} \ensuremath{-} 1); 272 \ensuremath{\Rightarrow} \textit{n} = 17
\item “At least one”: 1 \ensuremath{-} P(none); shared birthday among \textit{k}: 1 \ensuremath{-} {\ensuremath{\dfrac{365\, \cdot \, 364\, \cdot \, …\, \cdot \, (366\, -\, {k})}{365^{{k}}}}}
\end{itemize}\par{\color{black!72}Maps from |\textit{M}| = \textit{m} to |\textit{N}| = \textit{n} elements: all \textit{n}\textsuperscript{\textit{m}}; injective {\ensuremath{\dfrac{{n}!}{({n}\, -\, {m})!}}}; bijective \textit{n}! (\textit{m} = \textit{n}); onto $\sum_{{j}=0}^{{n}} $ (\ensuremath{-}1)\textsuperscript{\textit{j}}C(\textit{n}, \textit{j})(\textit{n} \ensuremath{-} \textit{j})\textsuperscript{\textit{m}}.}\par}
{\sheethead{Induction}\begin{enumerate}
\item Base: check A(\textit{n}\textsubscript{0}) for the first \textit{n}.
\item Step: assume A(\textit{n}) (hypothesis), prove A(\textit{n} + 1).
\item Conclude A(\textit{n}) for all \textit{n} \ensuremath{\geq} \textit{n}\textsubscript{0}.
\end{enumerate}\par{\color{black!72}Sums: S\textsubscript{\textit{n}+1} = S\textsubscript{\textit{n}} + (new term), insert the hypothesis, factor. Divisibility: 5\textsuperscript{\textit{n}+1} \ensuremath{-} 1 = 5(5\textsuperscript{\textit{n}} \ensuremath{-} 1) + 4. Inequality: Bernoulli (1 + \textit{x})\textsuperscript{\textit{n}} \ensuremath{\geq} 1 + \textit{n}\textit{x} for \textit{x} \ensuremath{\geq} \ensuremath{-}1. Two-step recurrences need two base cases.}\par}
{\sheethead{Maps}\begin{itemize}
\item \textit{f}: \textit{X} \ensuremath{\to} \textit{Y}, \textit{x} \ensuremath{\mapsto} \textit{f}(\textit{x}): domain \textit{X}, codomain \textit{Y}, one output per input
\item Image \textit{f}(\textit{A}) = \{\textit{f}(\textit{x}) : \textit{x} \ensuremath{\in} \textit{A}\}; Im \textit{f} = \textit{f}(\textit{X})
\item Preimage \textit{f}\textsuperscript{\ensuremath{-}1}(\textit{B}) = \{\textit{x} \ensuremath{\in} \textit{X} : \textit{f}(\textit{x}) \ensuremath{\in} \textit{B}\}: exists for every \textit{f}
\item Equal maps: same domain, codomain and rule
\end{itemize}}
{\sheethead{Injective, surjective, bijective}\begin{itemize}
\item \textbf{Injective}: \textit{f}(\textit{x}\textsubscript{1}) = \textit{f}(\textit{x}\textsubscript{2}) \ensuremath{\Rightarrow} \textit{x}\textsubscript{1} = \textit{x}\textsubscript{2} (each horizontal line hits at most once)
\item \textbf{Surjective} (onto): every \textit{y} \ensuremath{\in} \textit{Y} has \textit{x} with \textit{f}(\textit{x}) = \textit{y}, i.e. Im \textit{f} = \textit{Y}
\item \textbf{Bijective}: both (exactly one \textit{x} for each \textit{y})
\end{itemize}\par{\color{black!72}Prove injective: assume \textit{f}(\textit{x}\textsubscript{1}) = \textit{f}(\textit{x}\textsubscript{2}), deduce \textit{x}\textsubscript{1} = \textit{x}\textsubscript{2}. Disprove: two inputs, same output, e.g. (\textit{x} \ensuremath{-} 2)\textsuperscript{2} at 1 and 3. Prove onto: solve \textit{y} = \textit{f}(\textit{x}), check \textit{x} \ensuremath{\in} \textit{X}. Strictly monotone \ensuremath{\Rightarrow} injective (not conversely). Repair: restrict the domain to one side of a vertex, shrink the codomain to Im \textit{f}.}\par}
{\sheethead{Inverse and composition}\begin{itemize}
\item Bijective \textit{f} has \textit{f}\textsuperscript{\ensuremath{-}1}: \textit{Y} \ensuremath{\to} \textit{X}, \textit{f}\textsuperscript{\ensuremath{-}1}(\textit{y}) = \textit{x} \ensuremath{\Leftrightarrow} \textit{f}(\textit{x}) = \textit{y}; find it by solving \textit{y} = \textit{f}(\textit{x}) for \textit{x}
\item \textit{f}\textsuperscript{\ensuremath{-}1} \ensuremath{\circ} \textit{f} = id\textsubscript{\textit{X}}, \textit{f} \ensuremath{\circ} \textit{f}\textsuperscript{\ensuremath{-}1} = id\textsubscript{\textit{Y}}; graph mirrored in \textit{y} = \textit{x}
\item (\textit{g} \ensuremath{\circ} \textit{f})(\textit{x}) = \textit{g}(\textit{f}(\textit{x})): \textit{f} first; domain \{\textit{x} : \textit{f}(\textit{x}) \ensuremath{\in} D(\textit{g})\}
\item Associative, not commutative; (\textit{g} \ensuremath{\circ} \textit{f})\textsuperscript{\ensuremath{-}1} = \textit{f}\textsuperscript{\ensuremath{-}1} \ensuremath{\circ} \textit{g}\textsuperscript{\ensuremath{-}1}
\end{itemize}\par{\color{black!72}\textit{f}, \textit{g} injective \ensuremath{\Rightarrow} \textit{g} \ensuremath{\circ} \textit{f} injective; both onto \ensuremath{\Rightarrow} \textit{g} \ensuremath{\circ} \textit{f} onto; \textit{g} \ensuremath{\circ} \textit{f} injective \ensuremath{\Rightarrow} \textit{f} injective; \textit{g} \ensuremath{\circ} \textit{f} onto \ensuremath{\Rightarrow} \textit{g} onto. Range of \textit{g} \ensuremath{\circ} \textit{f}: find the range of \textit{f}, then apply \textit{g} to it.}\par}
{\sheethead{Polynomials and quadratics}\sheetf{\ensuremath{{a}{x}^{2}\, +\, {b}{x}\, +\, {c}\, =\, 0,\, {a}\, ≠\, 0}}\sheetf{\ensuremath{Δ\, =\, {b}^{2}\, -\, 4{a}{c},\, {x}\, =\, \dfrac{-{b}\, \pm \, \sqrt{Δ}}{2{a}}}}\begin{itemize}
\item Δ \ensuremath{>} 0: two real roots; Δ = 0: one double root; Δ \ensuremath{<} 0: none real
\item Vertex form: \textit{a}(\textit{x} + {\ensuremath{\dfrac{{b}}{2{a}}}})\textsuperscript{2} \ensuremath{-} {\ensuremath{\dfrac{Δ}{4{a}}}}
\item \textit{x}\textsubscript{0} is a root \ensuremath{\Leftrightarrow} (\textit{x} \ensuremath{-} \textit{x}\textsubscript{0}) divides P; degree \textit{n} \ensuremath{\Rightarrow} at most \textit{n} real roots
\item Biquadratic: \textit{u} = \textit{x}\textsuperscript{2} (or \textit{x}\textsuperscript{4}) turns it into a quadratic
\end{itemize}\par{\color{black!72}Parameter in front of \textit{x}\textsuperscript{2}: treat \textit{a} = 0 (a line) separately.}\par}
{\sheethead{exp, ln, powers}\begin{itemize}
\item e\textsuperscript{\textit{x}+\textit{y}} = e\textsuperscript{\textit{x}}e\textsuperscript{\textit{y}}, (e\textsuperscript{\textit{x}})\textsuperscript{\textit{y}} = e\textsuperscript{\textit{x}\textit{y}}, e\textsuperscript{0} = 1, e\textsuperscript{\textit{x}} \ensuremath{>} 0
\item exp: \ensuremath{\mathbb{R}} \ensuremath{\to} (0, \ensuremath{\infty}) strictly increasing, bijective; ln = exp\textsuperscript{\ensuremath{-}1}: (0, \ensuremath{\infty}) \ensuremath{\to} \ensuremath{\mathbb{R}}
\item ln(\textit{x}\textit{y}) = ln \textit{x} + ln \textit{y}, ln {\ensuremath{\dfrac{{x}}{{y}}}} = ln \textit{x} \ensuremath{-} ln \textit{y}, ln \textit{x}\textsuperscript{\textit{r}} = \textit{r} ln \textit{x}, ln 1 = 0, ln e = 1
\item \textit{a}\textsuperscript{\textit{x}} = e\textsuperscript{\textit{x} ln \textit{a}}; log\textsubscript{\textit{a}} \textit{x} = {\ensuremath{\dfrac{\ln \, {x}}{\ln \, {a}}}} (\textit{a} \ensuremath{>} 0, \textit{a} \ensuremath{\neq} 1)
\end{itemize}}
{\sheethead{Trigonometry}\par{\setlength{\tabcolsep}{2.4pt}\renewcommand{\arraystretch}{0.92}\begin{tabularx}{\linewidth}{@{}llllll>{\raggedright\arraybackslash}X@{}}
\textbf{\textit{x}} & \textbf{0} & \textbf{{\ensuremath{\dfrac{π}{6}}}} & \textbf{{\ensuremath{\dfrac{π}{4}}}} & \textbf{{\ensuremath{\dfrac{π}{3}}}} & \textbf{{\ensuremath{\dfrac{π}{2}}}} & \textbf{π} \\
sin & 0 & {\ensuremath{\dfrac{1}{2}}} & {\ensuremath{\dfrac{\sqrt{2}}{2}}} & {\ensuremath{\dfrac{\sqrt{3}}{2}}} & 1 & 0 \\
cos & 1 & {\ensuremath{\dfrac{\sqrt{3}}{2}}} & {\ensuremath{\dfrac{\sqrt{2}}{2}}} & {\ensuremath{\dfrac{1}{2}}} & 0 & \ensuremath{-}1 \\
tan & 0 & {\ensuremath{\dfrac{\sqrt{3}}{3}}} & 1 & \ensuremath{\sqrt{3}} & — & 0 \\
\end{tabularx}}\par\begin{itemize}
\item sin\textsuperscript{2} \textit{x} + cos\textsuperscript{2} \textit{x} = 1; sin(\ensuremath{-}\textit{x}) = \ensuremath{-}sin \textit{x}, cos(\ensuremath{-}\textit{x}) = cos \textit{x}; period 2π (tan: π)
\item sin(\textit{x} \ensuremath{\pm} \textit{y}) = sin \textit{x} cos \textit{y} \ensuremath{\pm} cos \textit{x} sin \textit{y}
\item cos(\textit{x} \ensuremath{\pm} \textit{y}) = cos \textit{x} cos \textit{y} \ensuremath{\mp} sin \textit{x} sin \textit{y}
\item tan(\textit{x} \ensuremath{\pm} \textit{y}) = {\ensuremath{\dfrac{\tan \, {x}\, \pm \, \tan \, {y}}{1\, \mp \, \tan \, {x}\, \tan \, {y}}}}
\item sin 2\textit{x} = 2 sin \textit{x} cos \textit{x}; cos 2\textit{x} = cos\textsuperscript{2} \textit{x} \ensuremath{-} sin\textsuperscript{2} \textit{x} = 2cos\textsuperscript{2} \textit{x} \ensuremath{-} 1
\item tan 2\textit{x} = {\ensuremath{\dfrac{2\, \tan \, {x}}{1\, -\, \tan ^{2}\, {x}}}}; cos 3\textit{x} = 4cos\textsuperscript{3} \textit{x} \ensuremath{-} 3cos \textit{x}
\item sin \textit{x} \ensuremath{-} sin \textit{y} = 2 cos {\ensuremath{\dfrac{{x}\, +\, {y}}{2}}} sin {\ensuremath{\dfrac{{x}\, -\, {y}}{2}}}
\item arcsin: [\ensuremath{-}1, 1] \ensuremath{\to} [\ensuremath{-}{\ensuremath{\dfrac{π}{2}}}, {\ensuremath{\dfrac{π}{2}}}]; arccos: [\ensuremath{-}1, 1] \ensuremath{\to} [0, π]; arctan: \ensuremath{\mathbb{R}} \ensuremath{\to} (\ensuremath{-}{\ensuremath{\dfrac{π}{2}}}, {\ensuremath{\dfrac{π}{2}}})
\item cos(arcsin \textit{x}) = \ensuremath{\sqrt{1\, -\, {x}^{2}}}; tan {\ensuremath{\dfrac{π}{12}}} = 2 \ensuremath{-} \ensuremath{\sqrt{3}}; cos {\ensuremath{\dfrac{π}{5}}} = {\ensuremath{\dfrac{1\, +\, \sqrt{5}}{4}}}
\end{itemize}}
{\sheethead{Finite decimals, computer numbers}\begin{itemize}
\item D = \{{\ensuremath{\dfrac{{m}}{10^{{n}}}}} : \textit{m} \ensuremath{\in} \ensuremath{\mathbb{Z}}, \textit{n} \ensuremath{\in} \ensuremath{\mathbb{N}}\textsubscript{0}\}: closed under + and \ensuremath{\cdot}; {\ensuremath{\dfrac{1}{3}}} \ensuremath{\notin} D
\item Truncation \textit{x}\textsubscript{\textit{n}} = {\ensuremath{\dfrac{{m}}{10^{{n}}}}}, \textit{m} = largest integer \ensuremath{\leq} 10\textsuperscript{\textit{n}}\textit{x}: 0 \ensuremath{\leq} \textit{x} \ensuremath{-} \textit{x}\textsubscript{\textit{n}} \ensuremath{<} 10\textsuperscript{\ensuremath{-}\textit{n}}
\item Floating point is not associative: (10\textsuperscript{16} \ensuremath{-} 10\textsuperscript{16}) + 1 = 1, but 10\textsuperscript{16} + (\ensuremath{-}10\textsuperscript{16} + 1) gives 0
\end{itemize}}
\end{multicols*}
\newpage
{\sffamily\bfseries\color{FAUblue}\small Side 2 \ensuremath{\cdot} Self-Study tasks, worked}\hfill{\fontsize{5}{6}\selectfont\color{black!65}Appendix B: Exam sheet, Exercises 1–4 · Mathematics for Engineers I \ensuremath{\cdot} Module 27651 \ensuremath{\cdot} Appendix B}\par\vspace{0.6pt}{\color{FAUblue}\hrule height 0.5pt}\vspace{1.2pt}
{\fontsize{4.6}{5.4}\selectfont\color{black!60}Personal exam sheet based on the MEI lecture notes and Exercise Sheets 1–4, in my own words; check what your exam allows before you take it in. Not an official FAU publication. Made by \textbf{Thanate Amornjiraporn (Focus)}.\par}\vspace{1pt}
\begin{multicols*}{4}\fontsize{8.80}{10.30}\selectfont\raggedright
{\sheethead{SS 1.4a \ensuremath{\cdot} statement or not?}\begin{itemize}
\item “If my father has a sister, …”: a statement (true or false, even if you do not know which)
\item “The formula is beautiful”: no definite truth value, so not a statement
\end{itemize}}
{\sheethead{SS 1.4c \ensuremath{\cdot} conditionals with false parts}\par C, D, P: cats, dogs, parrots have wings (all F).\par\begin{itemize}
\item cats only if dogs: C \ensuremath{\Rightarrow} D
\item parrots only if dogs: P \ensuremath{\Rightarrow} D
\item parrots if dogs: D \ensuremath{\Rightarrow} P
\item (C \ensuremath{\lor} D) \ensuremath{\Rightarrow} (C \ensuremath{\land} D)
\end{itemize}\par{\color{black!72}All four are true: each premise is false, and F \ensuremath{\Rightarrow} anything is T.}\par}
{\sheethead{SS 1.4e \ensuremath{\cdot} if, only if, unless}\par B: mirror breaks, G: glass, W: high on the wall.\par\begin{itemize}
\item breaks if glass: G \ensuremath{\Rightarrow} B
\item breaks only if glass: B \ensuremath{\Rightarrow} G
\item unless glass, not break: \ensuremath{\neg}G \ensuremath{\Rightarrow} \ensuremath{\neg}B \ensuremath{\Leftrightarrow} B \ensuremath{\Rightarrow} G (same as the last)
\item high, but not break unless glass: W \ensuremath{\land} (B \ensuremath{\Rightarrow} G)
\end{itemize}}
{\sheethead{SS 1.4d \ensuremath{\cdot} not a tautology}\par (\textit{A} \ensuremath{\Rightarrow} \textit{B}) \ensuremath{\Rightarrow} (\textit{B} \ensuremath{\Rightarrow} \textit{A}) with \textit{A} = F, \textit{B} = T: T \ensuremath{\Rightarrow} F = F. One false row is enough.\par}
{\sheethead{SS 1.4b \ensuremath{\cdot} De Morgan by truth table}\par{\setlength{\tabcolsep}{2.4pt}\renewcommand{\arraystretch}{0.92}\begin{tabularx}{\linewidth}{@{}ll>{\raggedright\arraybackslash}X@{}}
\textbf{\textit{A} \textit{B}} & \textbf{\ensuremath{\neg}(\textit{A} \ensuremath{\land} \textit{B})} & \textbf{\ensuremath{\neg}\textit{A} \ensuremath{\lor} \ensuremath{\neg}\textit{B}} \\
T T & F & F \\
T F & T & T \\
F T & T & T \\
F F & T & T \\
\end{tabularx}}\par\par{\color{black!72}Columns equal in every row \ensuremath{\Rightarrow} equivalent. The distributive law needs 8 rows.}\par}
{\sheethead{SS 1.5a \ensuremath{\cdot} intervals and points}\par \textit{A} = (\ensuremath{-}2, 1], \textit{B} = (\ensuremath{-}1, 1), \textit{C} = \{\ensuremath{-}2, 0, 1\} (roots of \textit{x}(\textit{x} + 2)(\textit{x} \ensuremath{-} 1)).\par\begin{itemize}
\item \textit{A} \ensuremath{\cap} \textit{B} = (\ensuremath{-}1, 1); \textit{A} \ensuremath{\cap} \textit{B} \ensuremath{\cap} \textit{C} = \{0\}
\item \textit{A} \ensuremath{\cap} (\textit{B} \ensuremath{\cup} \textit{C}) = (\ensuremath{-}1, 1]
\item (\textit{A} \ensuremath{\cup} \textit{B})\textsuperscript{c} = (\ensuremath{-}\ensuremath{\infty}, \ensuremath{-}2] \ensuremath{\cup} (1, \ensuremath{\infty})
\item (\textit{A} \textbackslash{} \textit{B}) \textbackslash{} \textit{C} = (\ensuremath{-}2, \ensuremath{-}1]
\end{itemize}}
{\sheethead{SS 1.5b \ensuremath{\cdot} power sets}\par \ensuremath{\mathcal{P}}(\{a, b, c\}) = \{\ensuremath{\emptyset}, \{a\}, \{b\}, \{c\}, \{a, b\}, \{a, c\}, \{b, c\}, \{a, b, c\}\}: 2\textsuperscript{3} = 8. \{1, 2\} \ensuremath{\times} \{1\} = \{(1, 1), (2, 1)\}, so \ensuremath{\mathcal{P}} has 4 members and \{1, 2\} \ensuremath{\notin} \ensuremath{\mathcal{P}}.\par}
{\sheethead{SS 1.5c \ensuremath{\cdot} a set identity}\sheetf{\ensuremath{{x}\, \in \, {A}\, \setminus \, ({B}\, \setminus \, {C})\, \Leftrightarrow \, {x}\, \in \, {A}\, \land \, \neg ({x}\, \in \, {B}\, \land \, {x}\, \notin \, {C})}}\sheetf{\ensuremath{\Leftrightarrow \, {x}\, \in \, {A}\, \land \, ({x}\, \notin \, {B}\, \lor \, {x}\, \in \, {C})\, \Leftrightarrow \, ({x}\, \in \, {A}\, \setminus \, {B})\, \lor \, ({x}\, \in \, {A}\, \cap \, {C})}}\par{\color{black!72}So \textit{A} \textbackslash{} (\textit{B} \textbackslash{} \textit{C}) = (\textit{A} \textbackslash{} \textit{B}) \ensuremath{\cup} (\textit{A} \ensuremath{\cap} \textit{C}): De Morgan, then distribute.}\par}
{\sheethead{SS 1.6 \ensuremath{\cdot} quantifiers in \ensuremath{\mathbb{N}}, \ensuremath{\mathbb{Z}}, \ensuremath{\mathbb{R}}\textsuperscript{+}}\par{\setlength{\tabcolsep}{2.4pt}\renewcommand{\arraystretch}{0.92}\begin{tabularx}{\linewidth}{@{}lll>{\raggedright\arraybackslash}X@{}}
\textbf{Claim} & \textbf{\ensuremath{\mathbb{N}}} & \textbf{\ensuremath{\mathbb{Z}}} & \textbf{\ensuremath{\mathbb{R}}\textsuperscript{+}} \\
\ensuremath{\exists}\textit{x} \ensuremath{\forall}\textit{y}: \textit{x} \ensuremath{<} \textit{y} & F & F & F \\
\ensuremath{\forall}\textit{y} \ensuremath{\exists}\textit{x}: \textit{x} \ensuremath{<} \textit{y} & F & T & T \\
\ensuremath{\exists}\textit{x} \ensuremath{\forall}\textit{y}: \textit{x} + \textit{y} = 1 & F & F & F \\
\ensuremath{\forall}\textit{y} \ensuremath{\exists}\textit{x}: \textit{x} + \textit{y} = 1 & F & T & F \\
\end{tabularx}}\par\par{\color{black!72}Row 1 fails at \textit{y} = \textit{x} (\textit{x} \ensuremath{<} \textit{x} is false). Row 2: \ensuremath{\mathbb{N}} has a least element 1; in \ensuremath{\mathbb{Z}} take \textit{x} = \textit{y} \ensuremath{-} 1, in \ensuremath{\mathbb{R}}\textsuperscript{+} \textit{x} = {\ensuremath{\dfrac{{y}}{2}}}. Row 3: \textit{y} = 1 forces \textit{x} = 0, \textit{y} = 2 forces \textit{x} = \ensuremath{-}1. Row 4: \textit{x} = 1 \ensuremath{-} \textit{y} stays in \ensuremath{\mathbb{Z}} only.}\par\par{\setlength{\tabcolsep}{2.4pt}\renewcommand{\arraystretch}{0.92}\begin{tabularx}{\linewidth}{@{}>{\raggedright\arraybackslash}X@{}}
\textbf{Negations} \\
\ensuremath{\forall}\textit{x} \ensuremath{\exists}\textit{y}: \textit{x} \ensuremath{\geq} \textit{y}; \ensuremath{\exists}\textit{y} \ensuremath{\forall}\textit{x}: \textit{x} \ensuremath{\geq} \textit{y} \\
\ensuremath{\forall}\textit{x} \ensuremath{\exists}\textit{y}: \textit{x} + \textit{y} \ensuremath{\neq} 1; \ensuremath{\exists}\textit{y} \ensuremath{\forall}\textit{x}: \textit{x} + \textit{y} \ensuremath{\neq} 1 \\
\end{tabularx}}\par}
{\sheethead{SS 2.5 \ensuremath{\cdot} regions and inequalities}\begin{itemize}
\item |\textit{x}| + |\textit{y} \ensuremath{-} 2| \ensuremath{<} 5: inside of the diamond with corners (\ensuremath{\pm}5, 2), (0, 7), (0, \ensuremath{-}3); edge excluded
\item |\textit{x} + 3| + |\textit{x} \ensuremath{-} 3| \ensuremath{>} 8: on [\ensuremath{-}3, 3] the sum is 6 (no); \textit{x} \ensuremath{>} 3: 2\textit{x} \ensuremath{>} 8 \ensuremath{\Rightarrow} \textit{x} \ensuremath{>} 4; \textit{x} \ensuremath{<} \ensuremath{-}3: \ensuremath{-}2\textit{x} \ensuremath{>} 8 \ensuremath{\Rightarrow} \textit{x} \ensuremath{<} \ensuremath{-}4. Answer (\ensuremath{-}\ensuremath{\infty}, \ensuremath{-}4) \ensuremath{\cup} (4, \ensuremath{\infty})
\item \ensuremath{\sqrt{({x}\, -\, 4)({x}\, -\, 1)}} \ensuremath{\leq} 2 \ensuremath{-} \textit{x}: radicand \ensuremath{\geq} 0 \ensuremath{\Leftrightarrow} \textit{x} \ensuremath{\leq} 1 or \textit{x} \ensuremath{\geq} 4; right side \ensuremath{\geq} 0 \ensuremath{\Leftrightarrow} \textit{x} \ensuremath{\leq} 2; so \textit{x} \ensuremath{\leq} 1. Square: \textit{x}\textsuperscript{2} \ensuremath{-} 5\textit{x} + 4 \ensuremath{\leq} \textit{x}\textsuperscript{2} \ensuremath{-} 4\textit{x} + 4 \ensuremath{\Leftrightarrow} \textit{x} \ensuremath{\geq} 0. Answer [0, 1]
\end{itemize}}
{\sheethead{SS 2.6 \ensuremath{\cdot} products and complements}\begin{itemize}
\item \textit{A} \ensuremath{\subset} \textit{C}, \textit{B} \ensuremath{\subset} \textit{D} \ensuremath{\Rightarrow} \textit{A} \ensuremath{\times} \textit{B} \ensuremath{\subset} \textit{C} \ensuremath{\times} \textit{D}: take (\textit{a}, \textit{b}), then \textit{a} \ensuremath{\in} \textit{C}, \textit{b} \ensuremath{\in} \textit{D}
\item Converse false: \textit{A} = \{1\}, \textit{B} = \ensuremath{\emptyset} gives \ensuremath{\emptyset} \ensuremath{\subset} anything; true if \textit{A}, \textit{B} \ensuremath{\neq} \ensuremath{\emptyset}
\item \textit{A} = [1, 3), \textit{B} = (2, 5]: \textit{A}\textsuperscript{c} = (\ensuremath{-}\ensuremath{\infty}, 1) \ensuremath{\cup} [3, \ensuremath{\infty}); \textit{A} \ensuremath{\cap} \textit{B} = (2, 3); \textit{A} \textbackslash{} \textit{B} = [1, 2]
\item ([1, 3] \ensuremath{\times} [1, 2])\textsuperscript{c}: \textit{x} \ensuremath{<} 1 or \textit{x} \ensuremath{>} 3 or \textit{y} \ensuremath{<} 1 or \textit{y} \ensuremath{>} 2, i.e. ((\ensuremath{-}\ensuremath{\infty}, 1) \ensuremath{\cup} (3, \ensuremath{\infty})) \ensuremath{\times} \ensuremath{\mathbb{R}} \ensuremath{\cup} \ensuremath{\mathbb{R}} \ensuremath{\times} ((\ensuremath{-}\ensuremath{\infty}, 1) \ensuremath{\cup} (2, \ensuremath{\infty}))
\item Complements: (\textit{A} \ensuremath{\cup} \textit{B})\textsuperscript{c} = (\ensuremath{-}\ensuremath{\infty}, 1) \ensuremath{\cup} (5, \ensuremath{\infty}); (\textit{A} \ensuremath{\cap} \textit{B})\textsuperscript{c} = (\ensuremath{-}\ensuremath{\infty}, 2] \ensuremath{\cup} [3, \ensuremath{\infty})
\end{itemize}}
{\sheethead{SS 2.7 \ensuremath{\cdot} sums, products, counting}\begin{itemize}
\item 10\textsuperscript{2} \ensuremath{-} {\ensuremath{\dfrac{1}{9^{2}}}} + 8\textsuperscript{2} \ensuremath{-} {\ensuremath{\dfrac{1}{7^{2}}}} + … + 2\textsuperscript{2} \ensuremath{-} {\ensuremath{\dfrac{1}{1^{2}}}} = $\sum_{{k}=0}^{4} $ ((10 \ensuremath{-} 2\textit{k})\textsuperscript{2} \ensuremath{-} {\ensuremath{\dfrac{1}{(9\, -\, 2{k})^{2}}}})
\item $\sum_{{k}=0}^{{n}} $ \textit{q}\textsuperscript{\textit{k}} \ensuremath{-} $\sum_{{k}=0}^{{n}+1} $ \textit{q}\textsuperscript{\textit{k}\ensuremath{-}1} = \ensuremath{-}{\ensuremath{\dfrac{1}{{q}}}} (all but \textit{q}\textsuperscript{\ensuremath{-}1} cancel)
\item 1 \ensuremath{-} {\ensuremath{\dfrac{2}{{j}({j}\, +\, 1)}}} = {\ensuremath{\dfrac{({j}\, -\, 1)({j}\, +\, 2)}{{j}({j}\, +\, 1)}}}, so $\prod_{{j}=2}^{{n}} $(…) = $\prod_{{j}=2}^{{n}} $ {\ensuremath{\dfrac{{j}\, -\, 1}{{j}}}} \ensuremath{\cdot} $\prod_{{j}=2}^{{n}} $ {\ensuremath{\dfrac{{j}\, +\, 2}{{j}\, +\, 1}}} = {\ensuremath{\dfrac{1}{{n}}}} \ensuremath{\cdot} {\ensuremath{\dfrac{{n}\, +\, 2}{3}}} = {\ensuremath{\dfrac{{n}\, +\, 2}{3{n}}}}
\item \ensuremath{\textstyle\sum} (\ensuremath{-}1)\textsuperscript{\textit{k}} C(\textit{n}, \textit{k}) = (1 \ensuremath{-} 1)\textsuperscript{\textit{n}} = 0
\item 6 marbles, \ensuremath{\leq} 2 striped, only 4 white: exactly 2 striped: C(12, 2) \ensuremath{\cdot} C(4, 4) = 66
\item 272 riddles = \textit{m}(\textit{m} \ensuremath{-} 1) \ensuremath{\Rightarrow} (\textit{m} \ensuremath{-} 17)(\textit{m} + 16) = 0 \ensuremath{\Rightarrow} 17 members
\end{itemize}}
{\sheethead{SS 3.5, 3.6 \ensuremath{\cdot} induction steps}\par Bases (\textit{n} = 1): 1 + \textit{x} = {\ensuremath{\dfrac{1\, -\, {x}^{2}}{1\, -\, {x}}}}; 1 = {\ensuremath{\dfrac{1\, \cdot \, 2}{2}}}; 1 = {\ensuremath{\dfrac{1\, \cdot \, 2\, \cdot \, 3}{6}}}; 1 = ({\ensuremath{\dfrac{1\, \cdot \, 2}{2}}})\textsuperscript{2}. Steps add the next term:\par\sheetf{\ensuremath{\sum_{{k}=0}^{{n}+1} \, {x}^{{k}}\, =\, \dfrac{1\, -\, {x}^{{n}+1}}{1\, -\, {x}}\, +\, {x}^{{n}+1}\, =\, \dfrac{1\, -\, {x}^{{n}+2}}{1\, -\, {x}}}}\sheetf{\ensuremath{\dfrac{{n}({n}\, +\, 1)}{2}\, +\, ({n}\, +\, 1)\, =\, \dfrac{({n}\, +\, 1)({n}\, +\, 2)}{2}}}\sheetf{\ensuremath{\dfrac{{n}({n}\, +\, 1)(2{n}\, +\, 1)}{6}\, +\, ({n}\, +\, 1)^{2}\, =\, \dfrac{({n}\, +\, 1)({n}\, +\, 2)(2{n}\, +\, 3)}{6},\, \text{as}\, 2{n}^{2}\, +\, 7{n}\, +\, 6\, =\, ({n}\, +\, 2)(2{n}\, +\, 3)}}\sheetf{\ensuremath{\dfrac{{n}^{2}({n}\, +\, 1)^{2}}{4}\, +\, ({n}\, +\, 1)^{3}\, =\, ({n}\, +\, 1)^{2}\left(\dfrac{{n}^{2}}{4}\, +\, {n}\, +\, 1\right)\, =\, \dfrac{({n}\, +\, 1)^{2}({n}\, +\, 2)^{2}}{4}}}}
{\sheethead{SS 3.7 \ensuremath{\cdot} classify and compose}\begin{itemize}
\item \textit{f}(\textit{x}) = (\textit{x} \ensuremath{-} 2)\textsuperscript{2} \ensuremath{-} 1: range [\ensuremath{-}1, \ensuremath{\infty}); not injective (\textit{f}(1) = \textit{f}(3) = 0), not onto \ensuremath{\mathbb{R}}. On [2, \ensuremath{\infty}) \ensuremath{\to} [\ensuremath{-}1, \ensuremath{\infty}): \textit{f}\textsuperscript{\ensuremath{-}1}(\textit{y}) = 2 + \ensuremath{\sqrt{{y}\, +\, 1}}
\item \textit{g}(\textit{x}) = \ensuremath{\sqrt{1\, -\, {x}}}, \textit{x} \ensuremath{\leq} 1: decreasing, injective, range [0, \ensuremath{\infty}); onto [0, \ensuremath{\infty}) with \textit{g}\textsuperscript{\ensuremath{-}1}(\textit{y}) = 1 \ensuremath{-} \textit{y}\textsuperscript{2}
\item \textit{h}(\textit{x}, \textit{y}, \textit{z}) = (\textit{x} + \textit{y}, \textit{y} + \textit{z}): onto (use (\textit{u}, 0, \textit{v})), not injective (add (\ensuremath{-}\textit{t}, \textit{t}, \ensuremath{-}\textit{t}))
\item \textit{f} \ensuremath{\circ} sin: range [0, 8]; \textit{g} \ensuremath{\circ} \textit{f} = \ensuremath{\sqrt{2\, -\, ({x}\, -\, 2)^{2}}} on [2 \ensuremath{-} \ensuremath{\sqrt{2}}, 2 + \ensuremath{\sqrt{2}}], range [0, \ensuremath{\sqrt{2}}]
\item \textit{g} \ensuremath{\circ} sin = \ensuremath{\sqrt{1\, -\, \sin \, {x}}}, range [0, \ensuremath{\sqrt{2}}]; \textit{g} \ensuremath{\circ} \textit{g} = \ensuremath{\sqrt{1\, -\, \sqrt{1\, -\, {x}}}} on [0, 1], range [0, 1]
\end{itemize}}
{\sheethead{SS 3.8 \ensuremath{\cdot} finite decimals}\begin{itemize}
\item {\ensuremath{\dfrac{{m}}{10^{{n}}}}} + {\ensuremath{\dfrac{{p}}{10^{{r}}}}} = {\ensuremath{\dfrac{{m}\, \cdot \, 10^{{N}-{n}}\, +\, {p}\, \cdot \, 10^{{N}-{r}}}{10^{{N}}}}} \ensuremath{\in} D with \textit{N} = max(\textit{n}, \textit{r}); product {\ensuremath{\dfrac{{m}{p}}{10^{{n}+{r}}}}} \ensuremath{\in} D
\item {\ensuremath{\dfrac{1}{3}}} = {\ensuremath{\dfrac{{m}}{10^{{n}}}}} \ensuremath{\Rightarrow} 10\textsuperscript{\textit{n}} = 3\textit{m}, but 10\textsuperscript{\textit{n}} leaves remainder 1 on division by 3
\item \textit{m} \ensuremath{\leq} 10\textsuperscript{\textit{n}}\textit{x} \ensuremath{<} \textit{m} + 1 \ensuremath{\Rightarrow} 0 \ensuremath{\leq} \textit{x} \ensuremath{-} {\ensuremath{\dfrac{{m}}{10^{{n}}}}} \ensuremath{<} 10\textsuperscript{\ensuremath{-}\textit{n}}
\item Floats: (10\textsuperscript{16} + (\ensuremath{-}10\textsuperscript{16})) + 1 = 1, but \ensuremath{-}10\textsuperscript{16} + 1 rounds to \ensuremath{-}10\textsuperscript{16} (neighbours are 2 apart there), so 10\textsuperscript{16} + (\ensuremath{-}10\textsuperscript{16} + 1) = 0
\end{itemize}}
{\sheethead{SS 3.9 \ensuremath{\cdot} Fibonacci in Pascal}\par \textit{D}\textsubscript{\textit{n}} = \ensuremath{\textstyle\sum} C(\textit{n} \ensuremath{-} \textit{k}, \textit{k}): 1, 1, 2, 3, 5, 8, 13, 21 = \textit{F}\textsubscript{\textit{n}+1}. Pascal inside the sum splits \textit{D}\textsubscript{\textit{n}+2} into \textit{D}\textsubscript{\textit{n}+1} + \textit{D}\textsubscript{\textit{n}}; two base cases \textit{D}\textsubscript{0} = \textit{D}\textsubscript{1} = 1.\par}
{\sheethead{SS 4.4 \ensuremath{\cdot} classify and compose}\begin{itemize}
\item \textit{f}(\textit{x}) = \ensuremath{-}\textit{x}\textsuperscript{2} + 6\textit{x} \ensuremath{-} 5 = 4 \ensuremath{-} (\textit{x} \ensuremath{-} 3)\textsuperscript{2}: range (\ensuremath{-}\ensuremath{\infty}, 4]; on [3, \ensuremath{\infty}): \textit{f}\textsuperscript{\ensuremath{-}1}(\textit{y}) = 3 + \ensuremath{\sqrt{4\, -\, {y}}}
\item \textit{g}(\textit{x}) = {\ensuremath{\dfrac{1}{1\, +\, {x}}}}, \textit{x} \ensuremath{\geq} 0: values (0, 1], decreasing; \textit{g}\textsuperscript{\ensuremath{-}1}(\textit{y}) = {\ensuremath{\dfrac{1}{{y}}}} \ensuremath{-} 1
\item \textit{h}(\textit{x}, \textit{y}) = (\textit{x}, \textit{x} + \textit{y}, \textit{y}): injective, image the plane \textit{b} = \textit{a} + \textit{c}, not onto \ensuremath{\mathbb{R}}\textsuperscript{3}
\item \textit{f} \ensuremath{\circ} exp: range (\ensuremath{-}\ensuremath{\infty}, 4]; \textit{g} \ensuremath{\circ} \textit{f} = {\ensuremath{\dfrac{1}{5\, -\, ({x}\, -\, 3)^{2}}}} on [1, 5], range [{\ensuremath{\dfrac{1}{5}}}, 1]
\item \textit{g} \ensuremath{\circ} exp = {\ensuremath{\dfrac{1}{1\, +\, \mathrm{e}^{{x}}}}}, range (0, 1); \textit{g} \ensuremath{\circ} \textit{g} = {\ensuremath{\dfrac{1\, +\, {x}}{2\, +\, {x}}}}, range [{\ensuremath{\dfrac{1}{2}}}, 1)
\end{itemize}}
{\sheethead{SS 4.5 \ensuremath{\cdot} composition facts}\begin{itemize}
\item \textit{g} \ensuremath{\circ} \textit{f} injective: \textit{f}(\textit{x}\textsubscript{1}) = \textit{f}(\textit{x}\textsubscript{2}) \ensuremath{\Rightarrow} \textit{g}(\textit{f}(\textit{x}\textsubscript{1})) = \textit{g}(\textit{f}(\textit{x}\textsubscript{2})) \ensuremath{\Rightarrow} \textit{x}\textsubscript{1} = \textit{x}\textsubscript{2}, so \textit{f} injective
\item \textit{g} \ensuremath{\circ} \textit{f} onto: \textit{z} = \textit{g}(\textit{f}(\textit{x})) gives \textit{y} = \textit{f}(\textit{x}) with \textit{g}(\textit{y}) = \textit{z}, so \textit{g} onto
\item \textit{f}, \textit{g} onto \ensuremath{\Rightarrow} \textit{g} \ensuremath{\circ} \textit{f} onto: for \textit{z} pick \textit{y} with \textit{g}(\textit{y}) = \textit{z}, then \textit{x} with \textit{f}(\textit{x}) = \textit{y}
\item \textit{g} \ensuremath{\circ} \textit{f} injective while \textit{g} is not: \textit{X} = \{a\}, \textit{f}(a) = 0, \textit{g}(0) = \textit{g}(1)
\end{itemize}}
{\sheethead{SS 4.6 \ensuremath{\cdot} roots; cos 3x}\begin{itemize}
\item 5\textit{x}\textsuperscript{4} + 20\textit{x}\textsuperscript{2} \ensuremath{-} 25: \textit{u} = \textit{x}\textsuperscript{2}, \textit{u}\textsuperscript{2} + 4\textit{u} \ensuremath{-} 5 = (\textit{u} \ensuremath{-} 1)(\textit{u} + 5) \ensuremath{\Rightarrow} 5(\textit{x} \ensuremath{-} 1)(\textit{x} + 1)(\textit{x}\textsuperscript{2} + 5); real roots \ensuremath{\pm}1, complex \ensuremath{\pm}i\ensuremath{\sqrt{5}}
\item \textit{a}\textit{x}\textsuperscript{2} + 2\textit{x} + 1: \textit{a} = 0 one root \ensuremath{-}{\ensuremath{\dfrac{1}{2}}}; Δ = 4(1 \ensuremath{-} \textit{a}): two for \textit{a} \ensuremath{<} 1 (\textit{a} \ensuremath{\neq} 0), double root \ensuremath{-}1 at \textit{a} = 1, none for \textit{a} \ensuremath{>} 1
\item cos 3\textit{x} = cos 2\textit{x} cos \textit{x} \ensuremath{-} sin 2\textit{x} sin \textit{x} = 4cos\textsuperscript{3} \textit{x} \ensuremath{-} 3 cos \textit{x} \ensuremath{\neq} 3 cos \textit{x} (try \textit{x} = 0)
\end{itemize}}
{\sheethead{SS 4.7 \ensuremath{\cdot} monotone and injective}\begin{itemize}
\item Strictly increasing or decreasing \ensuremath{\Rightarrow} injective: \textit{x} \ensuremath{\neq} \textit{y} gives \textit{f}(\textit{x}) \ensuremath{<} \textit{f}(\textit{y}) or \ensuremath{>}
\item Converse false: swap 0 and 1, fix the rest: bijective, not monotone
\item exp strictly increasing and onto (0, \ensuremath{\infty}) \ensuremath{\Rightarrow} bijective, inverse ln
\end{itemize}}
{\sheethead{SS 4.8 \ensuremath{\cdot} tangent formulas}\begin{itemize}
\item Divide the sin and cos addition formulas by cos \textit{x} cos \textit{y}: tan(\textit{x} \ensuremath{\pm} \textit{y}) = {\ensuremath{\dfrac{\tan \, {x}\, \pm \, \tan \, {y}}{1\, \mp \, \tan \, {x}\, \tan \, {y}}}}
\item tan 2\textit{x} = {\ensuremath{\dfrac{2\, \tan \, {x}}{1\, -\, \tan ^{2}\, {x}}}}
\item tan {\ensuremath{\dfrac{π}{12}}} = tan({\ensuremath{\dfrac{π}{3}}} \ensuremath{-} {\ensuremath{\dfrac{π}{4}}}) = {\ensuremath{\dfrac{\sqrt{3}\, -\, 1}{1\, +\, \sqrt{3}}}} = 2 \ensuremath{-} \ensuremath{\sqrt{3}}
\end{itemize}}
{\sheethead{Before you hand in}\begin{itemize}
\item Endpoints: open or closed? strict \ensuremath{<} or \ensuremath{\leq}? Check each end of an interval answer.
\item Domain first: radicand \ensuremath{\geq} 0, denominator \ensuremath{\neq} 0, log argument \ensuremath{>} 0; square only non-negative sides.
\item Multiplying an inequality by a negative number flips it.
\item A \ensuremath{\Rightarrow} B with A false is true; the converse is a different statement.
\item Negating: swap \ensuremath{\forall} \ensuremath{\leftrightarrow} \ensuremath{\exists} and negate only the final condition.
\item Maps: name domain and codomain before deciding onto; one counterexample disproves injective.
\item Induction: base case written out, hypothesis used in the step.
\item Counting: does order matter? are repeats allowed? then pick the formula.
\end{itemize}}
\end{multicols*}
\end{document}
